\section{有限维分布怎样描述一个过程}
\label{sec:06-Random-Processes/01-Foundations/02}

一个随机过程可能涉及无穷多个时间点。连续时间过程在任意有限区间内就有不可数个时刻，离散时间过程也有可数无限个。你没办法写出一张无穷维的联合概率表——那不是一个人类能操作或计算机能存储的对象。

但你有办法退而求其次。不需要一次性指定所有时刻的无穷维联合分布；只要保证：随便你挑出有限个时刻，我都能告诉你这些时刻上的联合分布是什么。

任取有限个时刻

\[
t_1, t_2, \ldots, t_n,
\]

考察随机向量

\[
(X_{t_1}, X_{t_2}, \ldots, X_{t_n})
\]

的联合分布。所有这样的（对所有有限 \(n\) 和所有时刻组合的）联合分布构成的集合，称为过程的\textbf{有限维分布族}（finite-dimensional distributions，常缩写为 fidis）。

这个想法很朴素：既然无穷维写不完，那我就保证所有有限维投影我都能回答。如果这个保证能被兑现，过程在概率层面就完整了——至少所有只涉及有限个时刻的事件，其概率都可以从有限维分布中计算。

\subsection{边缘分布远远不够——一维分布看不出依赖}

知道每个 \(X_t\) 各自的分布，你只拥有一维边缘分布。要判断时间依赖，你还需要二维联合分布 \((X_s, X_t)\)，以及三维、四维，直到任意有限维。

上一节构造的两个过程就是最直白的对照。过程 A：\(X_t = Z\)（随机水平线）。过程 B：每个整数时刻 \(Y_t\) 独立同分布 \(\mathcal N(0,1)\)。两个过程在每一个固定时刻 \(t\) 的分布完全相同：

\[
X_t \overset{d}{=} Y_t \sim \mathcal N(0,1).
\]

一维分布分不开它们。但看二维联合分布：

\begin{itemize}
\tightlist
\item
  过程 A：\((X_s, X_t)\) 的联合分布集中在直线 \(x = y\) 上——是一个退化的二元分布，相关系数为 1。
\item
  过程 B：当 \(s \neq t\) 时，\((Y_s, Y_t)\) 的联合分布是乘积分布——两个独立标准正态的联合，相关系数为 0。
\end{itemize}

二维分布立刻拉开了这两个过程。但二维分布本身也可能漏掉更高阶的关系。举个例子：你可以构造三个随机变量，两两独立但不相互独立（经典的``两两独立不蕴含相互独立''的例子）。四个、五个的情况类似——两两关系正常，不代表更高阶的联合结构没有隐藏的约束。

因此，完整指定一个过程的概率律需要所有有限维联合分布，而不仅仅是前几阶。

监控系统会直接暴露这种差别。假设 \(X_t\) 表示服务集群在第 \(t\) 分钟是否过载（1 = 过载，0 = 正常）。两个集群可能都满足 \(P(X_t = 1) = 0.1\)——单分钟过载比例完全相同，直方图对在一起分毫不差。但一个集群的过载事件近似逐分钟独立，另一个集群一旦进入过载就容易连续维持十几分钟。告警是否频繁连响、缓存是否需要更大容量，取决于 \((X_t, X_{t+1}, X_{t+2}, \ldots)\) 的联合分布——过载事件的持续长度和聚集模式，不是``百分之十的分钟过载''这个边缘比例能回答的。有限维分布记录的正是事件怎样沿时间聚集，而不仅仅是多少。

\subsection{有限维分布必须彼此一致——不能随意拼凑}

你不能任意为每组时刻指定一个联合分布。如果给 \((t_1, t_2)\) 指定了一个联合分布，又给 \((t_1, t_2, t_3)\) 指定了另一个，这两个指定必须相容——否则它们描述的就不可能是同一个过程。

相容条件叫一致性（consistency），两条：

\begin{itemize}
\tightlist
\item
  \textbf{排列对称}：交换时刻的顺序，联合分布只做相同的坐标置换。\((X_{t_1}, X_{t_2})\) 的联合分布和 \((X_{t_2}, X_{t_1})\) 的联合分布只是坐标位置交换——这不是两个独立的规定，而必须来自同一个底层概率律。
\item
  \textbf{边缘化相容}：从 \(n\) 维联合分布中边缘化掉某些时刻，应得到相应的低维联合分布。例如，若已指定 \((X_{t_1}, X_{t_2}, X_{t_3})\) 的三维联合分布，对其第三个坐标积分（求和）得到的二维边缘必须与先前指定的 \((X_{t_1}, X_{t_2})\) 的二维联合分布完全一致。
\end{itemize}

这两条不是数学上的吹毛求疵。如果你手工赋值的时候不顾一致性，你指定的``有限维分布''之间会互相矛盾——就像一个拼图的两块来自不同的盒子，拼不到一起。

Kolmogorov 扩张定理告诉你，只要一致性被满足，在足够好的状态空间上（通常的 \(\mathbb R\) 值或 \(\mathbb R^d\) 值过程都符合），存在一个随机过程，其有限维分布恰好是你指定的这一族。有限维分布因此不只是近似描述——它通过 Kolmogorov 扩张定理成为了构造随机过程的一条合法途径：你先在有限维层面把分布铺好，确保一致性，然后定理保证整个过程在无限维坐标空间上存在。

但这里有一个重要陷阱。有限维分布确定的是过程在坐标空间上的概率律——它告诉你在每个有限时刻子集上概率如何分配。但仅凭有限维分布，你无法读出所选版本的路径性质。同样的有限维分布族可以对应不同版本的路径——有的路径连续，有的路径疯狂跳跃。要把概率测度集中在连续函数空间或 càdlàg 路径空间（右连左极），需要额外的条件：矩条件、正则性定理（如 Kolmogorov 连续性定理），或者明确构造一个具有所需路径性质的过程版本。有限维分布定义了一个过程的``骨架''，但路径是骨架上长的肉——需要额外的工作来确保那些肉长在你想要的位置。

\subsection{Gaussian 过程：为什么只需均值和协方差}

对一般过程，有限维分布族是一大堆不同维度的联合分布，每一项都需要单独指定。但有一类过程让这件事极大地简化了。

如果对任意有限个时刻，随机向量

\[
(X_{t_1}, \ldots, X_{t_n})
\]

都服从多元 Gaussian 分布，就称 \(X_t\) 为 \textbf{Gaussian 过程}。

多元 Gaussian 的优雅之处在于：它完全由均值向量和协方差矩阵决定。对于整个过程，这意味着你只需要两个函数：

\[
m(t) = \E[X_t],
\qquad
K(s, t) = \Cov(X_s, X_t).
\]

只要 \(K\) 是合法的半正定核——即对任意有限个时刻 \(\{t_i\}\)，矩阵 \([K(t_i, t_j)]\) 半正定——就存在一个 Gaussian 过程以 \(m\) 为均值函数、\(K\) 为协方差函数。所有有限维分布自动由 \(m\) 和 \(K\) 通过多元 Gaussian 的公式生成。

这是一张巨大的简化票。对非 Gaussian 过程，均值和协方差只描述二阶结构（均值是位置，协方差是成对波动），不能确定所有有限维分布——你还需要更高阶的联合矩或者完整的联合分布函数。Gaussian 过程的二阶结构因多元 Gaussian 的封闭性而自动决定了所有阶的联合分布，不需要额外信息。

\subsection{独立增量与平稳增量——用局部规则构造全局结构}

直接给所有有限维分布赋值对多数过程来说太繁重。更常见的构造策略是：给定少数几条局部生成规则，让它们自动长出所有有限维联合分布。

两条最常用的增量规则：

\noindent\textbf{独立增量}：对于不重叠的时间区间，增量相互独立。正式地说，若区间 \((s_1, t_1], (s_2, t_2], \ldots, (s_k, t_k]\) 两两不重叠，则

\[
X_{t_1} - X_{s_1},\; X_{t_2} - X_{s_2},\; \ldots,\; X_{t_k} - X_{s_k}
\]

相互独立。这等于说过程在不相交的时间段上``忘记''自己之前做了什么——每个新区间的变化量与历史无关。

\noindent\textbf{平稳增量}：增量的分布只依赖区间长度，不依赖区间在时间轴上的绝对位置：

\[
X_{t+h} - X_{s+h} \overset{d}{=} X_t - X_s.
\]

平移整个区间，增量的分布不变。

Brownian motion 和 Poisson 过程都同时满足独立增量和平稳增量，但增量的分布完全不同：Brownian 的增量是 Gaussian，Poisson 的增量是 Poisson。平稳增量也不等于过程本身平稳——注意条件只约束增量，不约束水平值。标准 Brownian motion 的方差 \(\Var(B_t) = t\) 随 \(t\) 线性增长，所以 \(B_t\) 本身不平稳；下一篇会以随机游走为例给出可手动验证的完全相同的论证。

\subsection{从转移机制构造有限维分布——离散 Markov 链的示范}

离散状态 Markov 链是最直接的构造范例。不需要单独写出每组时刻的联合概率，只需要指定两个东西：

\begin{enumerate}
\tightlist
\item
  初始分布：\(P(X_0 = x_0)\)。
\item
  一步转移概率：\(P(X_{k} = x_k \mid X_{k-1} = x_{k-1})\)。
\end{enumerate}

这两个组件通过 Markov 性——给定现在，未来与过去条件独立——自动生成所有有限维联合分布：

\[
P(X_0 = x_0, X_1 = x_1, \ldots, X_n = x_n)
= P(X_0 = x_0) \prod_{k=1}^{n} P(X_k = x_k \mid X_{k-1} = x_{k-1}).
\]

每一步的条件概率只往前看一步，乘积链式法则负责把所有步骤串起来。在连续状态空间中，同一分解仍成立，但条件概率被替换为转移核（条件分布的积分表示）；若转移核有密度，记号上可以写成密度的乘积，处理方式类似。

Poisson 过程走的是增量路线：给定初始值 \(N(0) = 0\)、独立平稳增量，以及

\[
N(t) - N(s) \sim \operatorname{Poisson}(\lambda(t-s)),
\]

所有有限维分布就可以通过不相交区间增量的独立乘积性质构造出来。

这些构造策略的共同精神是：用少量的局部规则——一步转移概率，或者区间增量的分布——生成一个庞大的联合概率体系。有限维分布是你检查构造是否正确的检验标准，但构造本身不需要你手动写出每一个维度的联合分布。
